\begin{abstract}
Mathematical reasoning remains a formidable obstacle for language models due to its intricacy and formal structure. This work presents DeepSeekMath~7B, an extension of DeepSeek-Coder-Base-v1.5 7B that undergoes continued pre-training using 120B tokens related to mathematics, derived from Common Crawl, as well as additional data comprising natural language and source code. Achieving a noteworthy 51.7\% on the advanced MATH benchmark without external inference mechanisms or ensembling methods, DeepSeekMath~7B approaches the level of performance seen in Gemini-Ultra and GPT-4. When employing self-consistency over 64 samples, the score further increases to 60.9\% on the MATH benchmark. Two main innovations underpin the mathematical reasoning ability of DeepSeekMath~: Firstly, the model benefits from a carefully designed data selection pipeline that fully leverages extensive public web resources. Secondly, we propose Group Relative Policy Optimization (GRPO), a modified version of Proximal Policy Optimization (PPO) that strengthens mathematical reasoning and offers improvements in PPO's memory efficiency.
\end{abstract}

\section{Introduction}

The emergence of large language models (LLMs) has transformed mathematical reasoning methodologies within artificial intelligence, leading to considerable progress in areas such as quantitative reasoning benchmarks \citep{MATH} and geometry-focused assessments \citep{trinh2024solving}. In addition, LLMs now offer valuable support for humans addressing challenging mathematical questions \citep{tao}. Nonetheless, leading systems like GPT-4 \citep{gpt4} and Gemini-Ultra \citep{gemini} remain unavailable to the public, leaving open-source alternatives lagging significantly in terms of performance.

This work presents DeepSeekMath, a specialized language model tailored for mathematics that exhibits a marked advancement over existing open-source alternatives and nearly matches the mathematical proficiency of GPT-4 when evaluated on academic standards. The foundation of this model is the DeepSeekMath Corpus, an extensive, high-quality dataset containing 120B math tokens assembled for pre-training. To build this corpus, we apply a fastText-based classifier \citep{joulin2016fasttext} to the Common Crawl (CC) dataset. Initially, the classifier is developed with positive samples from OpenWebMath \citep{openwebmath}, while negative samples are drawn from a broad array of unrelated web pages. Using this initial model, we extract additional positive examples from CC, which are subsequently verified and refined by human annotators. The updated, curated data is then used to further train the classifier, enhancing its reliability. The quality of this large-scale corpus is validated by the performance of our base model, DeepSeekMath-Base 7B, which achieves 64.2\% on GSM8K \citep{gsm8k} and 36.2\% on the MATH competition dataset \citep{MATH}, surpassing the Minerva 540B model \citep{minerva}. Furthermore, since the DeepSeekMath Corpus includes multiple languages, we observe performance gains on Chinese mathematical benchmarks \citep{wei2023cmath,agieval}. We consider our methodology for mathematical data curation a valuable resource for the research community, acknowledging considerable opportunities for further advancements.

DeepSeekMath-Base is constructed by initializing with DeepSeek-Coder-Base-v1.5 7B \citep{deepseek-coder}, since we find that models with code-focused pre-training serve as a stronger foundation than those trained solely as general LLMs. Moreover, we observe that mathematical training extends benefits beyond mathematical competence, enhancing performance on general reasoning assessments such as MMLU \citep{mmlu} and BBH \citep{bbh}. This suggests that math-specific training confers broader reasoning gains to the model.

Following pre-training, we further fine-tune DeepSeekMath-Base on mathematical instruction datasets incorporating chain-of-thought \citep{cot}, program-of-thought \citep{pot,pal}, and tool-integrated reasoning \citep{tora} exemplars. This process yields DeepSeekMath-Instruct 7B, which surpasses all other 7B models and rivals open-source instruction-tuned models as large as 70B parameters.

Additionally, we propose Group Relative Policy Optimization (GRPO), a new reinforcement learning algorithm based on Proximal Policy Optimization (PPO) \citep{schulman2017proximal}. Unlike traditional PPO, GRPO discards the critic network and derives baselines from group-level scoring, leading to significantly reduced computational demands during training. Utilizing only a subset of English instruction-tuning data, GRPO achieves considerable gains over the already strong DeepSeekMath-Instruct across both domain-specific (GSM8K: 82.9\% $\rightarrow$ 88.2\%, MATH: 46.8\% $\rightarrow$ 51.7\%) and cross-domain mathematical benchmarks (such as CMATH: 84.6\% $\rightarrow$ 88.8\%) in the reinforcement learning phase. We also provide a unified framework that contextualizes methods such as Rejection Sampling Fine-Tuning (RFT) \citep{yuan2023scaling}, Direct Preference Optimization (DPO) \citep{dpo}, PPO, and GRPO, demonstrating how they represent direct or simplified instances of reinforcement learning. Our extensive experimental analysis—including studies on online vs. offline optimization, process vs. outcome-based supervision, and single-turn vs. iterative RL—further clarifies the critical components underpinning this framework. Finally, we elucidate how our RL methodology enhances instruction-tuned model performance and outline avenues for developing more effective RL strategies based on this unifying perspective.

\subsection{Contributions}
This work presents advancements in large-scale mathematical pre-training and delivers an in-depth investigation into the application of reinforcement learning.

\textbf{Math Pre-Training at Scale}

\begin{itemize}[topsep=0pt]
    \item We demonstrate that Common Crawl, an open-access dataset, harbors significant resources for mathematical machine learning. Through the development of a carefully engineered data filtering and selection process, we assemble the DeepSeekMath Corpus, consisting of 120B math-specific tokens extracted from the web. This dataset substantially surpasses the volume of mathematical web data utilized by Minerva \citep{minerva} by nearly sevenfold, and is about nine times larger than the latest OpenWebMath \citep{openwebmath} dataset.

    \item Our DeepSeekMath-Base 7B model, pre-trained on this specialized data, delivers performance on par with Minerva 540B \citep{minerva}, suggesting that scaling up model parameters alone does not guarantee enhanced mathematical reasoning. Models with fewer parameters, if trained on curated, high-quality data, can also reach competitive performance levels.

    \item Experimental results from our math training indicate that pre-training on code prior to engaging in mathematical tasks notably enhances a model’s capacity to address mathematical questions, irrespective of tool use. This contributes to the ongoing debate regarding the impact of code training on reasoning, supporting the hypothesis that such training can benefit mathematical reasoning abilities.

    \item Our results further reveal that although incorporating arXiv paper data is a standard practice in math-oriented training pipelines, doing so yields no measurable performance gains across all mathematical evaluation benchmarks employed in this study.
\end{itemize}

\textbf{Exploration and Analysis of Reinforcement Learning}

\begin{itemize}[topsep=0pt]
    \item We present Group Relative Policy Optimization (GRPO), a reinforcement learning technique that is both resource-efficient and effective. GRPO eliminates the need for a critic network by relying on group score-based baseline estimation, thereby substantially decreasing the computational demands compared to Proximal Policy Optimization (PPO).
    \item Our results show that applying GRPO substantially improves the performance of our instruction-tuned model, DeepSeekMath-Instruct, relying solely on instruction-tuning data. We also note performance improvements in out-of-domain scenarios throughout the reinforcement learning phase.
    \item We offer a cohesive framework to interpret a range of methodologies, including RFT, DPO, PPO, and GRPO. Through a comprehensive series of experiments—including online vs. offline optimization, supervision by outcome versus by process, and single-step vs. iterative reinforcement learning—we examine in depth the fundamental properties of this unified approach.
    \item Building on this conceptual framework, we analyze the underlying factors that drive the effectiveness of reinforcement learning and outline promising future avenues to realize more powerful reinforcement learning for LLMs.
\end{itemize}

\subsection{Summary of Evaluations and Metrics}

\begin{itemize}[topsep=0pt]
    \item \textbf{English and Chinese Mathematical Reasoning}: We rigorously evaluate our models on a suite of English and Chinese mathematics benchmarks spanning elementary to university-level questions. The English test sets include GSM8K \citep{gsm8k}, MATH \citep{MATH}, SAT \citep{llemma}, OCW Courses \citep{minerva}, and MMLU-STEM \citep{mmlu}; for Chinese, we use MGSM-zh \citep{mgsm}, CMATH \citep{wei2023cmath}, Gaokao-MathCloze \citep{agieval}, and Gaokao-MathQA \citep{agieval}. We measure each model's capacity to independently construct textual mathematical solutions and its competence at solving problems programmatically with Python.

    DeepSeekMath-Base achieves performance on par with the proprietary Minerva 540B \citep{minerva} model on English benchmarks and outperforms all publicly available base models, including Mistral 7B \citep{mistral} and Llemma-34B \citep{llemma}, regardless of prior math-focused pre-training. Particularly on Chinese benchmarks, DeepSeekMath-Base demonstrates marked superiority, a result attributable to our deliberate inclusion of high-quality data beyond English sources—unlike prior works \citep{minerva,llemma} that restrict pre-training to English materials. Incorporating instruction tuning and reinforcement learning into our models, the enhanced DeepSeekMath-Instruct and DeepSeekMath-RL attain a milestone in the open-source field by exceeding 50\% accuracy on the competition-grade MATH dataset for the first time.
\end{itemize}

\item \textbf{Formal Mathematics}: DeepSeekMath-Base is assessed on the informal-to-formal theorem proving task from \citep{dsp_proof} using miniF2F \citep{minif2f} as the evaluation set, with Isabelle \citep{isabelle} serving as the proof assistant. The results indicate that DeepSeekMath-Base attains robust few-shot autoformalization capabilities.

\item \textbf{Natural Language Understanding, Reasoning, and Code}: To comprehensively assess the model’s proficiency in general language understanding, reasoning, and programming, we benchmark DeepSeekMath-Base using the Massive Multitask Language Understanding (MMLU) dataset \citep{mmlu}, which consists of 57 diverse multiple-choice tasks, and BIG-Bench Hard (BBH) \citep{bbh}, which contains 23 demanding tasks predominantly involving multi-step reasoning. Additionally, we evaluate the model on HumanEval \citep{codex} and MBPP \citep{mbpp}, established benchmarks for code generation models. The results highlight the positive impact of math-centric pre-training on language understanding and reasoning abilities.

\section{Math Pre-Training}

\subsection{Data Collection and Decontamination} Here, we detail the methodology for assembling the DeepSeekMath Corpus utilizing data sourced from Common Crawl. As illustrated in Figure \ref{fig:data_collect}, our approach uses an iterative process for curating a large-scale mathematics dataset, initiated with a carefully selected seed set (for instance, a small but meticulously curated set of math datasets). Notably, this methodology is adaptable and can be employed for data collection in various other fields, such as programming.

Our process begins by adopting OpenWebMath \citep{openwebmath}—a dataset of curated mathematical web documents—as the seed collection. Leveraging this seed, we train a fastText model \citep{joulin2016fasttext} to retrieve additional web pages similar in style and content to those found in OpenWebMath. For model training, we sample 500,000 positive instances from the seed and pair them with 500,000 negative samples randomly drawn from Common Crawl. The open-source library\footnote{\url{https://fasttext.cc}} is used for model training, with hyperparameters set to a 256-dimensional vector space, learning rate of 0.1, a maximum word n-gram length of 3, minimum word occurrence threshold of 3, and 3 training epochs. To effectively shrink the scale of Common Crawl data, we implement both URL-based deduplication and near-duplicate detection, narrowing the dataset to 40B HTML web pages. Next, the fastText model is employed to identify and retrieve mathematical documents from the deduplicated set. We filter out low-quality pages by ranking all retrieved entries according to their fastText-assigned scores and select only those at the top of the ranking. The precise volume retained is empirically validated with pre-training experiments using the top 40B, 80B, 120B, and 160B tokens, and in this initial round, we select the top 40B tokens for further processing.

Following the initial round of data acquisition, a considerable portion of mathematical web content remains uncollected. This is primarily due to the limited diversity of the positive samples used to train the fastText model. To address this, we expand the pool of initial data sources by identifying further mathematics-related web domains, thereby enhancing the representativeness of the seed corpus and refining the fastText classifier. Concretely, we segment the entire Common Crawl dataset into non-overlapping domains, each defined by web pages that share a common base URL. For every domain, we determine the fraction of web pages that were included during the first collection phase. Those domains in which more than 10\% of web pages are collected are labeled as math-oriented (e.g., \url{mathoverflow.net}). 

We then manually review and label the URLs in these domains that are associated with mathematical material (such as \url{mathoverflow.net/questions}). Any web content corresponding to these annotated URLs, but not yet captured, is incorporated into the seed corpus. This process supplies a broader range of positive instances, which, in turn, allows for a more effective fastText model, improving its capacity to retrieve additional mathematical data in subsequent rounds. After four cycles of data collection, our resulting corpus comprises 35.5M math web pages, totaling 120B tokens. By the fourth cycle, we observe that almost 98\% of pages have already been obtained by the third cycle, prompting us to stop further data gathering.

To mitigate the risk of contaminating evaluation benchmarks, we apply filtering strategies consistent with \cite{deepseek-coder}. Specifically, we exclude any web pages containing content from major English math benchmarks (e.g., GSM8K \citep{gsm8k}, MATH \citep{MATH}) and Chinese math benchmarks (e.g., CMATH \citep{wei2023cmath}, AGIEval \citep{agieval}). Our filtering methodology dictates that any segment of text from our math corpus which includes a 10-gram sequence precisely matching a substring from any evaluation benchmark is removed. For cases where benchmark texts are under 10-grams but at least 3-grams long, we also employ exact matching to filter potentially contaminated content.

\subsection{Validating the Quality of the DeepSeekMath~Corpus}

To assess the utility of the DeepSeekMath~Corpus in comparison to other recently released math-focused corpora, we conduct pre-training experiments utilizing the following resources:

\begin{itemize}[topsep=0pt]
    \item \textbf{MathPile} \citep{mathpile}: a composite dataset (8.9B tokens) compiled from sources including textbooks, Wikipedia, ProofWiki, CommonCrawl, StackExchange, and arXiv, with the largest contribution (over 85\%) coming from arXiv;
    \item \textbf{OpenWebMath} \citep{openwebmath}: mathematical data extracted from CommonCrawl, comprising 13.6B tokens;
    \item \textbf{Proof-Pile-2} \citep{llemma}: an aggregation of OpenWebMath, AlgebraicStack (10.3B tokens of math code), and arXiv publications (28.0B tokens). In our experiments with Proof-Pile-2, we maintain the 2:4:1 ratio for arXiv:Web:Code as prescribed by \cite{llemma}.
\end{itemize}

\subsubsection{Training Setting}
\label{sec:quality-policy}
We adapt a math-specific training regime to a 1.3B-parameter general-purpose language model, which is architecturally aligned with DeepSeek LLMs \citep{deepseek-llm}, referenced here as DeepSeek-LLM 1.3B. Each mathematical dataset is utilized independently to train a distinct model over 150B tokens. The training pipeline relies on the HAI-LLM \citep{haillm} framework, notable for its efficiency and minimal resource overhead. Consistent with DeepSeek LLMs' established methodology, we use the AdamW optimizer \citep{adamW} with $\beta_1=0.9$, $\beta_2=0.95$, and $\mathrm{weight\_decay}=0.1$. The learning rate follows a multi-phase schedule: after 2,000 warmup steps, it climbs to its maximum, then reduces to 31.6\% at the 80\% mark of total training, and further drops to 10.0\% of the peak at the 90\% point. The learning rate peaks at 5.3e-4. Each training run employs a batch size of 4M tokens with a 4K context window.

\subsubsection{Evaluation Results}

\textbf{DeepSeekMath~Corpus demonstrates superior quality, multilingual breadth, and leads in scale among mathematical datasets.}

\begin{itemize}[topsep=0pt]
    \item \textbf{High-quality}:
    We assess downstream mathematical problem-solving by applying few-shot chain-of-thought prompting \cite{cot} across eight benchmarks. The model trained with DeepSeekMath~Corpus consistently achieves the highest results, as detailed in Table \ref{tab:corpora_comparison}. Additionally, Figure \ref{fig:corpora_comparisons} highlights that, even after only 50B tokens (matching one full Proof-Pile-2 epoch), training with DeepSeekMath yields superior accuracy, underscoring its higher overall data quality compared to other corpora.
    \item \textbf{Multilingual}:
    Data within the DeepSeekMath~Corpus spans several languages, with English and Chinese as the predominant contributors. Table \ref{tab:corpora_comparison} illustrates notable advancements in both English and Chinese mathematical reasoning following training with DeepSeekMath~Corpus. In comparison, available mathematical datasets are overwhelmingly English-based and deliver limited, sometimes detrimental, effects on Chinese-language mathematical tasks.
    \item \textbf{Large-scale}:
    DeepSeekMath~Corpus significantly surpasses the size of prior mathematical corpora. Figure \ref{fig:corpora_comparisons} reveals that DeepSeek-LLM 1.3B, when trained on this corpus, exhibits rapid learning progress and sustained performance gains. Baseline datasets, by contrast, are far smaller and must be repeatedly cycled during training, causing model improvements to quickly plateau.
\end{itemize}

\subsection{Training and Evaluating DeepSeekMath-Base 7B}

This section presents DeepSeekMath-Base 7B, a foundational model engineered for advanced reasoning capabilities, particularly within the domain of mathematics. The initialization of our model utilizes DeepSeek-Coder-Base-v1.5 7B \citep{deepseek-coder}, and the training process encompasses a total of 500B tokens. The composition of the training data is distributed as follows: 56\% derives from the DeepSeekMath Corpus, 4\% originates from AlgebraicStack, 10\% comes from arXiv, 20\% consists of Github code, while the remaining 10\% is made up of natural language material sourced from Common Crawl in both English and Chinese. The training procedure adheres primarily to the strategy described in Section \ref{sec:quality-policy}, with two exceptions: the learning rate peaks at 4.2e-4, and the batch size is set at 10 million tokens.

We perform an extensive evaluation of the mathematical proficiency of DeepSeekMath-Base 7B. This evaluation emphasizes its aptitude for independently generating complete mathematical solutions, its competency in resolving problems utilizing auxiliary tools, and its capacity for formal theorem verification. In addition to its mathematical assessment, we provide an overview of the model's general abilities, encompassing natural language comprehension, logical reasoning, and coding proficiency.

\paragraph{Mathematical Problem Solving with Step-by-Step Reasoning}

To measure DeepSeekMath-Base's skill in addressing mathematical questions, we employ few-shot chain-of-thought prompting \citep{cot} across eight distinct English and Chinese benchmarks. These datasets span tasks requiring quantitative reasoning, such as GSM8K \citep{gsm8k}, MATH \citep{MATH}, and CMATH \citep{wei2023cmath}, as well as multiple-choice evaluations like MMLU-STEM \citep{mmlu} and Gaokao-MathQA \citep{agieval}. Collectively, these benchmarks represent a wide array of mathematical disciplines, ranging from primary to collegiate levels of complexity.

As presented in Table \ref{tab:base_math_cot}, among all publicly available base models—spanning established architectures like Mistral 7B \citep{mistral} and the newer, mathematics-specialized Llemma 34B \citep{llemma} trained on Proof-Pile-2 \citep{llemma}—DeepSeekMath-Base 7B achieves the highest results across eight benchmarks. Particularly noteworthy is its performance on the challenging MATH dataset, where DeepSeekMath-Base surpasses all other open-source base models by more than 10\% in absolute terms, and even exceeds the results of the much larger, closed-source Minerva 540B model \citep{minerva}—an architecture 77 times the size, based on PaLM \citep{palm} and enhanced through further mathematical data.

\paragraph{Mathematical Problem Solving with Tool Use} For evaluating mathematical reasoning with programmatic tools, we conduct few-shot program-of-thought prompting \citep{pot,pal} on the GSM8K and MATH benchmarks. In this setup, models are instructed to tackle each task by crafting a Python script, with access to computational libraries such as \textit{math} and \textit{sympy} to handle advanced calculations. The model’s answer is determined by running the generated code. According to Table \ref{tab:base_math_pot_proof}, DeepSeekMath-Base 7B establishes a new standard, surpassing the previous leading model, Llemma 34B.

\paragraph{Formal Mathematics}

Automating formal proofs not only improves precision and dependability in mathematics, but also enhances productivity—a trend attracting rising interest. We assess DeepSeekMath-Base 7B on informal-to-formal proof generation as described in \citep{dsp_proof}: given an informal statement, its formal version, and a corresponding informal proof, the model produces a formalized proof. Evaluation occurs on the miniF2F benchmark \citep{minif2f}, which encompasses formal mathematics at the Olympiad level. Using a few-shot prompting approach, we instruct the model to generate formal proofs in Isabelle for each given problem. Following \cite{dsp_proof}, we direct the model to produce high-level proof sketches, which are then completed using the automated theorem prover Sledgehammer \citep{sledgehammer}. Table \ref{tab:base_math_pot_proof} highlights DeepSeekMath-Base 7B’s competitive results in automated proof formalization.

\paragraph{Natural Language Understanding, Reasoning, and Code}

The evaluation of DeepSeekMath-Base 7B’s capabilities in natural language understanding uses MMLU \citep{mmlu}, in reasoning utilizes BBH \citep{bbh}, and in coding is conducted on HumanEval \citep{codex} and MBPP \citep{mbpp}. Data in Table \ref{tab:understanding_reasoning_code} reveals that DeepSeekMath-Base 7B delivers substantial improvements on both MMLU and BBH compared to its predecessor, DeepSeek-Coder-Base-v1.5 \citep{deepseek-coder}, indicating that mathematical training augments language comprehension and logical reasoning. Furthermore, by integrating code-related data during continued training, DeepSeekMath-Base 7B sustains the coding benchmark performance previously seen in DeepSeek-Coder-Base-v1.5. In summary, DeepSeekMath-Base 7B outperforms the general-purpose Mistral 7B \citep{mistral} across all three reasoning and coding assessments.

\section{Supervised Fine-Tuning}

\subsection{SFT Data Curation}

We develop a comprehensive instruction-tuning dataset for mathematics, comprising both English and Chinese questions from a variety of mathematical domains and encompassing a wide spectrum of difficulty levels. Each problem is paired with a corresponding solution formatted as chain-of-thought (CoT) \citep{cot}, program-of-thought (PoT) \citep{pot,pal}, or in a tool-augmented reasoning style \citep{tora}. The resulting dataset contains 776K examples for training.

\begin{itemize}[topsep=0pt]
    \item \textbf{English mathematical datasets}: Tool-augmented answers are annotated for GSM8K and MATH datasets. Additionally, a portion of the MathInstruct dataset \citep{MathInstruct} and the Lila-OOD training set \citep{lila}—where solutions use CoT or PoT methods—are included. The English collection spans a broad range of areas, including but not limited to algebra, probability, number theory, calculus, and geometry.
    \item \textbf{Chinese mathematical datasets}: We gather Chinese K-12 level problems, encompassing 76 different sub-topics such as linear equations. For each, solutions are provided both in CoT format and with tool-integrated reasoning.
\end{itemize}

\subsection{Training and Evaluating DeepSeekMath-Instruct 7B}

This section details the training protocol for DeepSeekMath-Instruct 7B, which is adapted from DeepSeekMath-Base using instruction tuning tailored for mathematics. For training, problem-solution examples are concatenated at random until the input sequence reaches up to 4K tokens. The model is fine-tuned for 500 training steps using a batch size of 256 and a constant learning rate of 5e-5.

The evaluation of model performance includes both settings: with and without tool assistance. Assessment is performed across four quantitative reasoning benchmarks, covering both English and Chinese tasks. Comparative results are reported alongside contemporaneous state-of-the-art systems:

\begin{itemize}[topsep=0pt]
    \item \textbf{Closed-source models} consist of: (1) the GPT family, with GPT-4 \citep{gpt4} and the GPT-4 Code Interpreter~\footnote{\url{https://openai.com/blog/chatgpt-plugins\#\#code-interpreter}} representing the strongest candidates, (2) Gemini Ultra and Gemini Pro \citep{gemini}, (3) Inflection-2 \citep{inflection-2}, (4) Grok-1~\footnote{\url{https://x.ai/model-card}},
    as well as models released more recently by Chinese technology companies such as (5) Baichuan-3~\footnote{\url{https://www.baichuan-ai.com}},
    and (6) the new GLM-4~\footnote{\url{https://open.bigmodel.cn/dev/api\#glm-4}}

    from the GLM series \citep{glm}.
\end{itemize}

These models are designed for broad applications and have typically been subjected to multiple alignment procedures.
\item \textbf{Open-source models} in this analysis encompass: (1) DeepSeek-LLM-Chat 67B \citep{deepseek-llm}, (2) Qwen 72B \citep{qwen}, (3) SeaLLM-v2 7B \citep{seallm}, and (4) ChatGLM3 6B \citep{chatglm3}, which are all general-purpose; in addition, we include models with explicit improvements in mathematical reasoning such as (5) InternLM2-Math 20B\footnote{\url{https://github.com/InternLM/InternLM-Math}}, an extension of InternLM2 trained further on mathematics before being subjected to instruction tuning, (6) Math-Shepherd-Mistral 7B, which augments Mistral 7B \citep{mistral} with PPO training \citep{schulman2017proximal} using a process-level supervised reward system, (7) the WizardMath series \citep{wizardmath}, which boosts Mistral 7B and Llama-2 70B \citep{llama2} in mathematical reasoning through evolve-instruct (a variant of instruction tuning utilizing AI-generated instructions) and PPO with a curriculum sourced primarily from GSM8K and MATH datasets, (8) MetaMath 70B \citep{metamath}, based on Llama-2 70B fine-tuned with a data-augmented GSM8K and MATH, (9) ToRA 34B \cite{tora}, which adapts CodeLlama 34B for tool-assisted mathematical problem solving, and (10) MAmmoTH 70B \citep{MathInstruct}, which is Llama-2 70B fine-tuned on MathInstruct via instruction tuning.
\end{itemize}

According to Table \ref{tab:sft_rl_math}, in the evaluation setting that prohibits external tool usage, DeepSeekMath-Instruct 7B achieves notable proficiency in multi-step reasoning tasks. Particularly on the competitive MATH benchmark, it outperforms all open-source models and most closed-source systems (including Inflection-2 and Gemini Pro) by a margin of at least 9 percentage points, maintaining superiority even relative to significantly larger architectures (such as Qwen 72B) and models explicitly optimized for mathematics via reinforcement learning (e.g., WizardMath-v1.1 7B). DeepSeekMath-Instruct's performance is on par with leading Chinese closed-source models such as GLM-4 and Baichuan-3 on MATH, though it continues to lag behind GPT-4 and Gemini Ultra.

When evaluated under conditions permitting both natural language reasoning and code-based tool use, DeepSeekMath-Instruct 7B approaches a 60\% accuracy rate on the MATH benchmark, thereby exceeding the performance of all previously available open-source alternatives. On other datasets, this model’s outcomes are comparable to DeepSeek-LLM-Chat 67B—the former state-of-the-art, which has ten times more parameters.

\section{Reinforcement Learning}

\subsection{Group Relative Policy Optimization} Reinforcement learning (RL) has been empirically validated as an effective method for further enhancing the mathematical reasoning capabilities of large language models (LLMs) subsequent to Supervised Fine-Tuning (SFT) \citep{wang2023math,wizardmath}. This section details Group Relative Policy Optimization (GRPO), our novel and resource-efficient RL technique.

\subsubsection{From PPO to GRPO} Proximal Policy Optimization (PPO) \citep{schulman2017proximal} is a widely adopted actor-critic reinforcement learning algorithm, especially popular for fine-tuning LLMs \citep{ouyang2022training}. Specifically, PPO seeks to optimize LLM parameters by maximizing the following surrogate objective:

\begin{equation}
\footnotesize
    \mathcal{J}_{PPO}(\theta) = \mathbb{E}{[q \sim P(Q), o \sim \pi_{\theta_{old}}(O|q)]} \frac{1}{|o|} \sum_{t=1}^{|o|} \min \left[ \frac{\pi_\theta(o_{t} | q, o_{<t})}{\pi_{\theta_{old}}(o_{t} | q, o_{<t})} A_{t}, \text{clip} \left( \frac{\pi_\theta(o_{t} | q, o_{<t})}{\pi_{\theta_{old}}(o_{t} | q, o_{

Here, $\pi_{\theta}$ and $\pi_{\theta_{old}}$ denote the updated and preceding policy models, respectively, while $q$ and $o$ represent questions and sampled outputs from both the question corpus and the former policy $\pi_{\theta_{old}}$. The hyper-parameter $\varepsilon$, integral to the Proximal Policy Optimization (PPO) algorithm, governs the clipping mechanism for training stability. $A_t$ refers to the advantage estimate, which is computed via Generalized Advantage Estimation (GAE) \citep{gae} utilizing rewards $\{r_{\ge t}\}$ in conjunction with a learned value function $V_{\psi}$. Consequently, PPO necessitates simultaneous training of a value function alongside the policy; to counteract excessive exploitation of the reward model, the typical strategy augments the per-token reward with a KL divergence penalty term relative to a reference policy at every token \citep{ouyang2022training}:

\begin{equation}
    r_{t} = r_\varphi(q, o_{\le t}) - \beta \log\frac{\pi_{\theta}(o_{t}|q, o_{<t})}{\pi_{ref}(o_{t}|q, o_{<t})},
\label{eq:PPO-reward}
\end{equation}

with $r_\varphi$ denoting the output of the reward model, $\pi_{ref}$ as the reference policy (commonly the initial SFT checkpoint), and $\beta$ acting as the coefficient scaling the KL regularization.

Training a value function for PPO, which is generally architecturally similar to the policy network, imposes significant demands in terms of both memory and computation. In the process of RL training, this value function is utilized as a baseline in advantage calculations to help reduce variance. Within the LLM domain, reward signals are often only provided at the final output token by the reward model, posing challenges for constructing a token-level accurate value function. To address these limitations, and as illustrated in Figure \ref{fig:grpo}, we introduce Group Relative Policy Optimization (GRPO). GRPO eliminates the necessity for learning an explicit value function, as required in PPO, by substituting the conventional baseline with the mean reward from a set of sampled outputs (for the same question). Specifically, for a given question $q$, GRPO samples a collection of outputs $\{o_1, o_2, \cdots, o_G\}$ using the previous policy $\pi_{\theta_{old}}$, and updates the current policy via maximization of the following formulation:

\begin{equation}
\footnotesize
\begin{split}
    \mathcal{J}_{GRPO}(\theta) &= \mathbb{E}{[q \sim P(Q), \{o_i\}_{i=1}^G \sim \pi_{\theta_{old}}(O|q)]}  \\
    & \frac{1}{G}\sum_{i=1}^G\frac{1}{|o_i|} \sum_{t=1}^{|o_i|} \left\{ \min \left[ \frac{\pi_\theta(o_{i,t} | q, o_{i,<t})}{\pi_{\theta_{old}}(o_{i,t} | q, o_{i,<t})} \hat{A}_{i,t}, \text{clip} \left( \frac{\pi_\theta(o_{i,t} | q, o_{i,<t})}{\pi_{\theta_{old}}(o_{i,t} | q, o_{i,<t})}, 1 - \varepsilon, 1 + \varepsilon \right)  \hat{A}_{i,t} \right] - \beta \mathbb{D}_{KL}\left[\pi_{\theta} || \pi_{ref}\right]\right\} ,
\end{split}
\label{eq:GRPO-obj}
\end{equation}

where both $\varepsilon$ and $\beta$ remain as hyper-parameters, and $\hat{A}_{i,t}$ is defined as the advantage, but calculated in relation to the other outputs within the sampled group, details of which will be elaborated in the following subsections. This group-based advantage formulation is closely aligned with the comparative methodology intrinsic to reward models, since such models are predominantly trained using datasets with intra-question output rankings. Importantly, GRPO differs from prior approaches by adding the KL regularization term directly as a penalty on the loss, rather than incorporating it into the reward, which simplifies the computation of $\hat{A}_{i,t}$. Furthermore, as opposed to the KL term utilized in (\ref{eq:PPO-reward}), the KL divergence in our approach is

\begin{algorithm}[t]
  \small
  \caption{Iterative Group Relative Policy Optimization}
  \textbf{Input} initial policy model $\pi_{\theta_{\text{init}}}$; reward models $r_\varphi$; task prompts $\mathcal{D}$; 
  hyperparameters $\varepsilon$, $\beta$, $\mu$
  \begin{algorithmic}[1]
    \State Initialize policy model: $\pi_\theta \leftarrow \pi_{\theta_{\text{init}}}$
    \For{iteration = 1, \dots, I}
       \State Set reference model $\pi_{ref} \leftarrow \pi_{\theta}$
      \For{step = 1, \dots, M}
      \State Select a batch $\mathcal{D}_b$ from $\mathcal{D}$
      \State Assign old policy model: $\pi_{\theta_{old}} \leftarrow \pi_{\theta}$ 
      \State For each prompt $q \in \mathcal{D}_b$, generate $G$ completions $\{o_i\}_{i=1}^G \sim \pi_{\theta_{old}} (\cdot \mid q) $
      \State Use $r_{\varphi}$ to assign rewards $\{r_i\}_{i=1}^{G}$ for each output $o_i$
      \State Estimate token-level advantages $\hat{A}_{i,t}$ for the $t$-th token of $o_i$ using group-based relative advantage calculation
      \For{GRPO iteration = 1, \dots, $\mu$}
        \State Update $\pi_{\theta}$ by maximizing the GRPO criterion (see Equation \ref{eq:GC-GRPO})
      \EndFor
    \EndFor \State Continually improve $r_\varphi$ through online training with replay.
    \EndFor 
  \end{algorithmic}
  \textbf{Output} $\pi_\theta$
  \label{alg:iter-grpo}
\end{algorithm}

\subsubsection{Outcome Supervision RL with GRPO} In this framework, each question $q$ is associated with a set of $G$ completions $\{o_1, o_2, \cdots, o_G\}$ produced by the previous policy $\pi_{\theta_{old}}$. The reward model assesses these completions, producing a set of scores $\mathbf{r}=\{r_1, r_2, \cdots, r_G\}$. These reward values are standardized by removing the group mean and scaling by the group standard deviation. Under outcome supervision, the normalized reward for each output $o_i$ serves as a uniform advantage across all of its tokens: $\hat{A}_{i, t} = \widetilde{r}_i = \frac{r_i- {\rm mean}(\mathbf{r})}{{\rm std}(\mathbf{r})}$. The policy parameters are then updated to maximize the target defined in equation (\ref{eq:GRPO-obj}).

\subsubsection{Process Supervision RL with GRPO} Since outcome-level feedback can be sparse or suboptimal for complex mathematical reasoning, following the methodology of \cite{wang2023math}, we consider process supervision, which supplies feedback after each intermediate reasoning step. More precisely, given a prompt $q$ and $G$ sampled completions $\{o_1, o_2, \cdots, o_G\}$, a process reward model scores the outputs at each step, resulting in reward sequences: $\mathbf{R} = \{ \{r_1^{index(1)},\cdots,r_1^{index(K_1)}\}, \cdots,  \{r_G^{index(1)},\cdots,r_G^{index(K_G)}\} \}$. Here, $index(j)$ refers to the position of the terminal token for the $j$-th step, and $K_i$ denotes the step count for output $o_i$. We standardize each step’s reward by subtracting the overall mean and dividing by the standard deviation: $\widetilde{r}_i^{index(j)} = \frac{r_i^{index(j)} - {\rm mean(\mathbf{R})}}{{\rm std(\mathbf{R})}}$.
Process supervision then computes the token-wise advantage by summing the normalized rewards from all subsequent reasoning steps: $\hat{A}_{i, t} = \sum_{index(j) \ge t} \widetilde{r}_i^{index(j)}$, and proceeds to optimize the policy model according to the objective in equation (\ref{eq:GRPO-obj}).

\subsubsection{Iterative RL with GRPO} As reinforcement learning training advances, the previous reward model may lose its effectiveness in providing adequate supervision for the evolving policy model. To address this, we investigate iterative RL using GRPO. As illustrated in Algorithm \ref{alg:iter-grpo}, the iterative GRPO approach constructs new reward model training datasets from samples generated by the current policy model. The reward model is continuously refined using a replay buffer that integrates 10\% of prior data. Next, the current policy model is designated as the new reference model, and further training proceeds by updating the policy model with feedback from the updated reward model.

\subsection{Training and Evaluating DeepSeekMath-RL}

Our RL experiments utilize DeepSeekMath-Instruct 7B as the foundation. The RL training set comprises chain-of-thought-formatted items from GSM8K and MATH present in the SFT corpus, totaling roughly 144K questions. We omit other SFT datasets to rigorously analyze RL’s effects on benchmarks that remain unseen during RL. Construction of the reward model’s training data adheres to the procedure outlined by \citep{wang2023math}. The initial reward model is built on DeepSeekMath-Base 7B, employing a 2e-5 learning rate. For GRPO training, the policy model uses a learning rate of 1e-6, with the KL coefficient fixed at 0.04. For each input, 64 responses are sampled. Sequence length and batch size are capped at 1024. The policy model undergoes a single update step per exploration round. Evaluation is carried out on the same benchmarks as DeepSeekMath-Instruct 7B. For DeepSeekMath-RL 7B, the GSM8K and MATH tasks, formulated in chain-of-thought style, are treated as in-domain; all other benchmark datasets are considered out-of-domain.

Table \ref{tab:sft_rl_math} presents results for both open-source and proprietary models on English and Chinese tasks, evaluating their chain-of-thought and tool-enhanced reasoning abilities. Our observations indicate the following: (1) When employing chain-of-thought reasoning, DeepSeekMath-RL 7B achieves accuracy rates of 88.2\% on GSM8K and 51.7\% on MATH. These outcomes exceed those of all other open-source models ranging from 7B to 70B parameters and outperform most closed-source counterparts as well. (2) Notably, DeepSeekMath-RL 7B’s instruction tuning was exclusively performed on chain-of-thought format data from GSM8K and MATH, with training initiated from the DeepSeekMath-Instruct 7B checkpoint. Despite the relatively narrow scope of its training corpus, DeepSeekMath-RL 7B surpasses DeepSeekMath-Instruct 7B on every evaluated metric, highlighting the significant advantages of reinforcement learning.

\section{Discussion}
In this section, we outline our principal observations from both pre-training and RL experimentation.

\subsection{Insights Gained from Pre-Training}

We begin by discussing our pre-training experiences. Unless explicitly stated otherwise, the following reflections are based on the experimental setup described in Section \ref{sec:quality-policy}. It should be mentioned that our use of the term DeepSeekMath Corpus within this context refers specifically to an 89-billion-token dataset derived from the second cycle of our data gathering.

\subsubsection{Code Training Enhances Mathematical Reasoning}

There exists a widely held but previously unproven belief that code-oriented pre-training may augment a model's reasoning skills. We seek to shed light on this question, specifically regarding mathematical problem-solving: our findings indicate that code-based training bolsters models’ capacity to perform mathematical reasoning, irrespective of the integration of external tools.

To systematically examine the influence of code pre-training on mathematical reasoning, we conducted analyses using both two-stage and single-stage training protocols:

\textbf{Two-Stage Training}

\begin{itemize}[topsep=0pt]
    \item \textbf{Code Training for 400B Tokens $\rightarrow$ Math Training for 150B Tokens}: The DeepSeek-LLM 1.3B model is first exposed to 400B tokens derived from code sources, followed by an additional 150B tokens focused on mathematical content;
    \item \textbf{General Training for 400B Tokens $\rightarrow$ Math Training for 150B Tokens}: As a baseline for comparison, we also conduct an experiment where the initial 400B tokens originate from a general large-scale corpus curated by DeepSeek-AI, rather than code, in order to explore whether code-based pretraining yields superior improvements in mathematical reasoning compared to a generic corpus.
\end{itemize}

\textbf{One-Stage Training}

\begin{itemize}[topsep=0pt]
    \item \textbf{Math Training for 150B Tokens}: The DeepSeek-LLM 1.3B model is exclusively trained on 150B math-focused tokens;
    \item \textbf{Training on a mixture of 400B Code Tokens and 150B Math Tokens}: Sequential math training after prior code training tends to impair the model’s performance on coding tasks. To address this, we explore whether co-training on a blend of both 400B code and 150B math tokens in a single stage can both support mathematical reasoning abilities and lessen catastrophic forgetting.
\end{itemize}

\paragraph{Results}
Tables \ref{tab:code-math} and \ref{tab:code-math-general-eval} summarize the model's downstream task results across these various training configurations.

Engaging in code pretraining enhances the model's capacity for program-based mathematical reasoning, regardless of whether a two-stage or a single-stage paradigm is used. According to Table \ref{tab:code-math}, in the two-stage process, code-only pretraining robustly boosts performance on tasks such as GSM8K and MATH when solved with Python, and this is further augmented by subsequent math-specific training. Furthermore, with one-stage joint training, simultaneous exposure to code and math tokens alleviates the catastrophic forgetting problem observed in sequential two-stage training, while improving both coding task accuracy (Table \ref{tab:code-math-general-eval}) and program-aided mathematical reasoning capabilities (Table \ref{tab:code-math}).

Moreover, code-centric pretraining bolsters mathematical reasoning even when external tools are not used. During two-stage training, beginning with code yields measurable gains and amplifies the effectiveness of the later math-focused stage, culminating in the highest observed results. Conversely, when code and math tokens are intermixed from the outset (one-stage), the model's tool-free mathematical reasoning is somewhat diminished. This may be attributed to the relatively modest parameter count of DeepSeek-LLM 1.3B, which may be insufficient to deeply encode both programming and mathematical paradigms at once.

\subsubsection{ArXiv Papers Appear Unhelpful for Advancing Mathematical Reasoning}

ArXiv papers are a frequent component of mathematical pre-training datasets \citep{minerva,gpt-f,llemma,mathpile}. Nevertheless, their actual effect on enhancing mathematical reasoning has not been thoroughly investigated. Contrary to possible expectations, our empirical results indicate that arXiv papers contribute little, if at all, to mathematical reasoning abilities. To assess this, we conducted experiments with models of varying capacities, specifically DeepSeek-LLM 1.3B and DeepSeek-Coder-Base-v1.5 7B \citep{deepseek-coder}, trained on arXiv-sourced datasets processed by different means:

\begin{itemize}[topsep=0pt]
    \item \textbf{MathPile} \citep{mathpile}: an 8.9-billion-token dataset, generated by applying heuristic cleaning and filtering, with scientific arXiv articles making up more than 85\%;
    \item \textbf{ArXiv-RedPajama} \citep{redpajama}: a 28-billion-token corpus including all arXiv LaTeX files, post removal of preambles, comments, macros, and bibliographies.
\end{itemize}

We separately pre-train DeepSeek-LLM 1.3B (150B tokens) and DeepSeek-Coder-Base-v1.5 7B (40B tokens) on each of these arXiv datasets. The outcomes suggest that pre-training exclusively on arXiv content does not benefit, and in some cases may even impair, mathematical reasoning skills. Across a range of mathematical assessments used in this work—ranging from quantitative problem sets like GSM8K and MATH (see Table \ref{tab:arxiv-cot}), to multiple-choice tests such as MMLU-STEM (Table \ref{tab:arxiv-cot}), as well as formal math benchmarks like miniF2F (Table \ref{tab:arxiv_atp})—neither model demonstrated measurable gains.

It is important, however, to acknowledge several caveats to this conclusion. Our evaluation does not address:

\begin{itemize}[topsep=0pt]
    \item Possible advantages of arXiv-based training for narrowly focused mathematical applications outside the current scope, such as transforming formal statements or proofs into informal versions (informalization);
    \item The influence of arXiv data when blended with other types of training data;
    \item Whether larger models might reveal different trends regarding arXiv papers’ impact.
\end{itemize}

Accordingly, these outstanding questions warrant further examination in future research.

\subsection{Insights on Reinforcement Learning}
\subsubsection{Toward a Unified Framework} This section aims to establish a unified framework for analyzing a variety of training paradigms, including SFT, RFT, DPO, PPO, and GRPO, as well as experimentally investigating the influences of components in this unified perspective. In general, the gradient of a training objective with respect to the parameters $\theta$ can be expressed as follows:

\begin{equation}
    \nabla_{\theta}\mathcal{J}_{\textcolor{red}{\mathcal{A}}}(\theta) = \mathbb{E}[\underbrace{(q,o) \sim \textcolor{red}{\mathcal{D}}}_{Data \ Source}]\left( \frac{1}{|o|} \sum_{t=1}^{|o|}  \underbrace{GC_{{\mathcal{A}}}(q, o, t, \textcolor{red}{\pi_{{rf}}})}_{Gradient \ Coefficient}  \nabla_{\theta}\log \pi_{\theta}(o_t | q, o_{<t})\right).
\label{eq:objective}
\end{equation}

This objective involves three principal elements: 1) \textit{Data Source $\mathcal{D}$}, which specifies the dataset for training; 2) \textit{Reward Function $\pi_{{rf}}$}, which supplies the reward signals used during optimization; and 3) \textit{Algorithm $\mathcal{A}$}, responsible for transforming both training data and reward signals into a gradient coefficient $GC$ that governs the extent of positive or negative updates for each example. Under this framework, we consider several prototypical approaches:

\begin{itemize}
    \item  \textbf{Supervised Fine-tuning (SFT)}: SFT involves adapting a pretrained model through further training on a curated dataset annotated by humans.
    \item \textbf{Rejection Sampling Fine-tuning (RFT)}: In RFT, the SFT model undergoes additional training using a subset of outputs generated by itself, filtered to ensure correct answers to SFT prompts.
    \item \textbf{Direct Preference Optimization (DPO)}: DPO enhances the SFT model by fine-tuning it with extra sampled responses and employing a pairwise preference-based loss.
    \item \textbf{Online Rejection Sampling Fine-tuning (Online RFT)}: Online RFT contrasts with RFT by continually updating the policy model, sampling augmented responses from the model as it evolves during training.
    \item \textbf{PPO/GRPO}: These approaches start from the SFT model as a base policy and apply reinforcement by sampling outputs from the current, continually updated policy.
\end{itemize}

The elements comprising these methods are detailed in Table \ref{tab:data_policy}. For a comprehensive exposition of the derivation, consult Appendix \ref{app:analysis-rl}.

\paragraph{Observation about Data Source} We categorize data sources into two primary types: online sampling and offline sampling. Online sampling involves collecting training samples from the ongoing, in-situ policy model during exploration, whereas offline sampling collects data using the pretrained SFT model. Methods such as RFT and DPO employ offline sampling, while Online RFT and GRPO utilize online sampling approaches.

Referencing Figure \ref{fig:rl-analysis}, we observe that Online RFT exhibits marked improvement over RFT across both benchmarks. At the initial training phase, the performance of Online RFT parallels that of RFT; however, as training progresses, Online RFT achieves a pronounced superiority, underscoring the strengths of online data generation. This observation can be rationalized by noting that, during early training, the actor model and the SFT model are still very similar, with their sampled outputs showing negligible variance. Conversely, as training advances, the data produced by the actor deviates more substantially, and sampling in real time increasingly confers a performance benefit.

\paragraph{Observation about Gradient Coefficient} In these algorithms, the conversion of reward signals into gradient coefficients serves to guide updates of model parameters. In our experimental framework, we distinguish between `Rule'-based and `Model'-based reward functions. The former classifies responses by assessing answer accuracy, whereas the latter utilizes a dedicated reward model—trained using rule-based evaluations—to score responses. Equations \ref{eq:GC-RFT} and \ref{eq:GC-GRPO} elucidate a fundamental distinction: GRPO dynamically modifies its gradient coefficient in accordance with the reward magnitude assigned by the reward model, enabling nuanced promotion and discouragement of responses based on the extent of their reward. In contrast, Online RFT treats all correct answers equivalently, reinforcing them uniformly and offering no punitive mechanism for incorrect outputs.

As illustrated in Figure \ref{fig:rl-analysis}, GRPO outperforms online RFT, emphasizing the effectiveness of adjusting positive and negative gradient coefficients. Moreover, GRPO+PS achieves better results than GRPO+OS, underscoring the advantage of employing fine-grained, step-aware gradient coefficients. We also investigate the use of iterative RL; in our experiments, we carry out two iterations. According to Figure \ref{fig:iter-rl}, iterative RL produces substantial gains in performance, particularly after the initial iteration.

\subsubsection{Why RL Works?} In this work, we employ reinforcement learning using a subset of instruction tuning data, resulting in considerable improvements over the instruction-tuned model. To shed further light on the mechanism behind RL’s effectiveness, we measure Pass@K and Maj@K accuracy for both the Instruct and RL models on two evaluation sets. Figure \ref{fig:maj-pass} reveals that while RL leads to better Maj@K outcomes, it does not enhance Pass@K. These observations suggest that RL mainly improves the robustness of the model’s output distribution, i.e., \textbf{the observed gains stem from promoting correct answers within TopK, rather than from enhancing core capabilities.} Likewise, \citep{wang2023making} identified a \textbf{misalignment problem} in reasoning tasks within SFT models, and demonstrated that the reasoning performance of SFT models could be elevated through several preference alignment approaches \citep{yuan2023rrhf,song2023preference,wang2023making}.

\subsubsection{How to Achieve More Effective RL?}
\label{sec:effec_rl}
We establish that RL is highly effective in mathematical reasoning scenarios. Furthermore, we introduce a unified framework that contextualizes various prominent training methods. Under this framework, all methods can be regarded as variants or simplifications of RL approaches. As encapsulated by Equation \ref{eq:objective}, the framework involves three principal elements: Data Source, Algorithm, and Reward Function. We propose potential future research avenues for each of these components.

\paragraph{Data Source} The data source serves as the foundational input for all training regimes. Within the RL setting, it specifically refers to unlabeled questions alongside candidate outputs generated by the policy model. Our experiments utilize only instruction-tuning stage questions and rely on a basic nucleus sampling method to generate outputs. This may partially account for the finding that our RL approach only raises Maj@K scores. Moving forward, we intend to test our RL strategy on out-of-distribution questions, together with \textbf{advanced sampling (decoding) strategies}, such as those leveraging tree-search algorithms \citep{yao2023tree}. In addition, the application of \textbf{efficient inference techniques} \citep{xia-etal-2023-speculative,leviathan2023fast,kwon2023efficient,xia2024unlocking}, which shape the exploration efficiency of policy models, remains a critical area for further investigation.

\paragraph{Algorithms} Algorithms utilize data and reward signals to determine the gradient coefficient used for adjusting model parameters. According to Equation \ref{eq:objective}, contemporary approaches universally \textbf{TRUST} the reward signal to guide increases or decreases in the conditional probability for particular tokens. Nonetheless, the reliability of reward signals cannot be guaranteed, especially in highly intricate scenarios. For instance, even the PRM800K datasets \citep{lightman2023let}, which are annotated by highly trained experts, still feature an error rate of about 20\% in their annotations\footnote{\url{https://github.com/openai/prm800k/issues/12\#issuecomment-1728491852}}. Therefore, it is crucial to investigate reinforcement learning methods that exhibit robustness in the presence of noisy reward signals. We anticipate that alignment strategies grounded in the \textbf{WEAK-TO-STRONG} paradigm \citep{burns2023weak} have the potential to fundamentally transform learning algorithms.

\paragraph{Reward Function} The reward function provides the foundational feedback required for training. In reinforcement learning, this function typically takes the form of a neural reward model. We identify three pivotal areas of focus for improving reward models: 1) \textbf{Strategies for improving the generalization capacity of the reward model.} A reward model needs to generalize well to cope with out-of-distribution queries and sophisticated decoding behaviors; failing this, RL may simply consolidate existing language model distributions rather than facilitating meaningful advancement; 2) \textbf{Capturing and expressing the uncertainty of reward models.} Quantifying uncertainty may establish an essential connection between limited reward models and weak-to-strong alignment algorithms; 3) \textbf{Methods for constructing process reward models with high efficiency and quality}, enabling them to deliver detailed feedback on reasoning processes \citep{lightman2023let,wang2023math}.

\section{Conclusion, Limitation, and Future Work}

We introduce DeepSeekMath, a model that achieves superior results among open-source alternatives on the competition-level MATH benchmark and exhibits performance close to that of proprietary models. DeepSeekMath~originates from DeepSeek-Coder-v1.5 7B, undergoing continual training with 500B tokens, including a substantial proportion—120B tokens—of mathematical content harvested from Common Crawl. Our comprehensive ablation analyses reveal that web data holds considerable promise as a rich source for mathematical datasets, whereas arXiv does not contribute as significantly as anticipated. We propose Group Relative Policy Optimization (GRPO), an adaptation of Proximal Policy Optimization (PPO), which enables enhanced mathematical reasoning while using less memory. Experimental outcomes confirm the efficacy of GRPO, even for DeepSeekMath-Instruct 7B when already performing strongly on benchmarks. Furthermore, we propose a unified framework to interpret a range of related methods and outline multiple promising avenues for advancing reinforcement learning.

Despite the strong results on quantitative reasoning benchmarks, DeepSeekMath~exhibits relatively lower proficiency in domains like geometry and theorem proving compared to closed-source models. For example, preliminary tests indicate the model's inability to effectively solve questions involving triangles and ellipses, suggesting a potential bias in data selection during pre-training and fine-tuning. Additionally, given its scale, DeepSeekMath~underperforms GPT-4 on few-shot tasks; unlike GPT-4, which shows clear performance gains from few-shot prompts, DeepSeekMath~delivers comparable results in both zero-shot and few-shot evaluations. Future efforts will focus on refining our engineered data selection methodology to assemble even higher quality pre-training corpora. We will also investigate new strategies (Section \ref{sec:effec_rl}) for optimizing reinforcement learning processes for large language models.

\end{document}